\usepackage[paper=b5paper]{geometry} %https://tex.stackexchange.com/questions/220139/b5paper-option-for-book-document-class-abnormal-results

\usepackage[utf8]{inputenc} % codificacio dels caracters
\usepackage[T1]{fontenc} % Nou esquema de codificacio (recomanat)
\usepackage{lmodern} % latin modern 

% \usepackage{ae,aecompl}
\usepackage[british]{babel} % localitzacio de l'estil
\usepackage[british]{isodate} % per les dates
\usepackage[hyphens]{url}
% \usepackage{amssymb,amsmath} % macros per matematiques
\usepackage{amsmath,amssymb} 
\usepackage[pdftex]{graphicx}	% support per figures
%\usepackage{epsfig}
\usepackage{lscape} % per poder posar figures i/o taules en horitzontal
\usepackage{color} % per poder fer servir colors
\usepackage{bibentry}

\usepackage{colortbl}


\nobibliography*
\usepackage[colorlinks,bookmarks,pdftex,hyperfigures,breaklinks,backref=page,
    pdfpagemode=UseOutlines,
    pdftitle=EnricGusoPhDThesis,
    pdfauthor=Enric~Guso,
    pdfsubject=PhD~Thesis,
    pdfkeywords=ecological validity~music information retrieval~deep learning
]{hyperref} % produeix enllacos clicables
\usepackage[authoryear]{natbib} % per poder fer citacions Author-any
\usepackage{fancyvrb}
\usepackage{units}
\usepackage[caption=false]{subfig}
%\usepackage{natbibspacing}
%\setlength{\bibspacing}{0.3\baselineskip}
% per les pagines a la bibliografia
% -> UN COMMENT A PARTIR D?AQUI
\renewcommand*{\backref}[1]{}
\renewcommand*{\backrefalt}[4]{%
    \ifcase #1 (Not cited.)%
    \or        (Cited on page~#2.)%
    \else      (Cited on pages~#2.)%
    \fi}

\usepackage{latexsym}
\usepackage{rotating}
\usepackage{bm}
%\usepackage{bibentry}\nobibliography*
\usepackage{longtable}
%\usepackage{verbatim}
% \newsubfloat{figure}
\usepackage{wrapfig}
% better tables
\usepackage[]{threeparttable}
\usepackage{booktabs}
\usepackage{arydshln}
\newcommand{\ra}[1]{\renewcommand{\arraystretch}{#1}}
\usepackage{array}
\newcolumntype{P}[1]{>{\raggedright\arraybackslash}p{#1}}
\newcolumntype{L}[1]{>{\centering\arraybackslash}m{#1}}
\newcolumntype{C}[1]{>{\centering\arraybackslash}p{#1}}
\newcolumntype{T}[1]{>{\raggedleft\arraybackslash}p{#1}}
\newcolumntype{H}{>{\setbox0=\hbox\bgroup}c<{\egroup}@{}}  %Hide column in a table
% for footnotes in tables
\usepackage{footnote}
\usepackage{multirow}
\usepackage{makecell}
\usepackage[toc]{glossaries}
\usepackage{doi}
\usepackage{rotating}
\usepackage[htt]{hyphenat}

\makenoidxglossaries
\newglossaryentry{latex}
{
    name=latex,
    description={LaTeX (short for Lamport TeX) is a document preparation system. The user has to think about only the content to put in the document and the software will take care of the formatting. }
}

\newglossaryentry{glsy}
{
    name=glossary,
    description={Acronyms and terms which are generally unknown or new to common readers.}
}

\newacronym{MSS}{MSS}{Music Source Separation}
\newacronym{IAM}{IAM}{Indian Art Music}
\newacronym{STFT}{STFT}{Short-Time Fourier Transfer}
\newacronym{MIR}{MIR}{Music Information Retrieval}
\newacronym{DNN}{DNN}{Deep Neural Network}
\newacronym{CNN}{CNN}{Convolutional Neural Network}
\newacronym{LSTM}{LSTM}{Long Short-Term Memory}
\newacronym{HAs}{HAs}{Hearing Aids}
\newacronym{SE}{SE}{Speech Enhancement}
\newacronym{RNN}{RNN}{Recurrent Neural Network}
\newacronym{RIR}{RIR}{Room Impulse Response}
\newacronym{GRU}{GRU}{Gated Recurrent Unit}
\newacronym{TCN}{TCN}{Temporal Convolutional Network}
\newacronym{GPU}{GPU}{Graphics Processing Unit}
\newacronym{PA}{PA}{Public Address}
\newacronym{GMACs}{GMACs}{Giga Multiply-Add Operations}
\newacronym{ASR}{ASR}{Automatic Speech Recognition}
\newacronym{VAE}{VAE}{Variational AutoEncoder}
\newacronym{SGD}{SGD}{Stochastic Gradient Descent}
\newacronym{NLP}{NLP}{Natural Language Processing}
\newacronym{SI-SDR}{SI-SDR}{Scale-Invariant Signal-to-Distortion Ratio}
\newacronym{SI-SDRi}{SI-SDRi}{Scale-Invariant Signal-to-Distortion Ratio Improvement}
\newacronym{SOTA}{SOTA}{State-of-the-Art}
\newacronym{GAN}{GAN}{Generative Adversarial Network}
\newacronym{GPT}{GPT}{Generative Pre-trained Transformer}
\newacronym{SNR}{SNR}{Signal to Noise Ratio}

\usepackage{ifthen} % per treure figures de apendix de la llista de figures, etc...

% poerque els footnotes no tornin a començar a cada chapter
\usepackage{chngcntr}
\counterwithout{footnote}{chapter}

\maxsecnumdepth{subsubsection} % Per defecte la classe 'memoir' no numera les subseccions. Que les numeri:
\setsecnumdepth{subsubsection}
\maxtocdepth{subsection} % Tampoc no inclou les subseccions a l'index. Que les inclogui:
\settocdepth{subsection}

\usepackage[font=small,labelfont=bf]{caption}

% palonso packages
\usepackage{soul}  % allows highlight
\usepackage{listings} % for code snippets
\lstset{
  language=Python,
  basicstyle=\ttfamily\small,
  numbers=left,
  numberstyle=\tiny,
  breaklines=true,
  keywordstyle=\color{blue},
  commentstyle=\color{gray}\itshape,
  stringstyle=\color{red!60!black},
  showstringspaces=false
}
\usepackage{enumitem} % numbered lists
\usepackage{subfiles} % including subfiles
\usepackage[x11names,table,dvipsnames]{xcolor} % http://ctan.org/pkg/xcolor
\usepackage{pifont} % commands for Pi fonts (Dingbats, Symbol, etc.)
\usepackage{wasysym}
\usepackage{fontawesome5}



% gcortes packages
\usepackage{bbm}
\usepackage{musicography}
\usepackage{subcaption}
\usepackage{cleveref}
\usepackage{rotating}
\usepackage{xargs} % for \newcommandx
\usepackage{longtable}
\usepackage{adjustbox}  % adjust table width
\usepackage{float} % Use H if don't want image/table repositioning
\usepackage[most]{tcolorbox}  % Colored tags
\usepackage{epigraph}
\usepackage{algorithm}
\usepackage{algpseudocode}

\tcbuselibrary{skins,breakable}
\newtcolorbox{mybox}[2][]{breakable, skin=enhancedmiddle jigsaw, parbox=false,
boxrule=0mm, leftrule=2mm, boxsep=0mm, arc=5mm, outer arc=5mm, attach title to upper,
after title={.\ }, coltitle=black, colback=white, colframe=gray!80, title={#2},
fonttitle=\bfseries, before=\vspace{0.15cm}, #1}
% Define a new tcolorbox environment with a thinner left rule and a fancy color
\newtcolorbox{myboxhighlight}[2][]{breakable, skin=enhancedmiddle jigsaw, parbox=false,
boxrule=0mm, leftrule=1mm, boxsep=0mm, arc=5mm, outer arc=5mm, attach title to upper,
after title={.\ }, coltitle=black, colback=white, colframe=teal!60!black, title={#2}, fonttitle=\bfseries,
before=\vspace{0.15cm}, #1}


\newlist{todolist}{itemize}{2}
\setlist[todolist]{label=$\square$}
\usepackage[colorinlistoftodos,prependcaption,textsize=scriptsize]{todonotes}